%! Function Operations — Unit 1, Lesson 1.2 — Lesson Plan
% Section order per LESSON_SHAPE.md §5. A *New* lesson — a foundations ramp with no CED topic.
% 1.1 taught notation, evaluating, domain/range and piecewise rules; this lesson runs two
% functions at once. The crux is the domain of a combination read off the SIMPLIFIED formula
% instead of off the ORIGINAL functions: (f/g) with f(x) = x^2 - 9 and g(x) = x - 3 simplifies
% to x + 3, which looks defined everywhere, but x = 3 is still excluded. 1.3 chains the two
% machines instead of combining them.
\documentclass[10pt]{article}
\usepackage{precalculus-article}
\usepackage{precalculus-boxes}
\usepackage{graphicx}
\graphicspath{{images/}}
\newcommand{\TallMath}[1]{$\displaystyle #1\rule[-1.4em]{0pt}{3.2em}$}
\newcommand{\CourseName}{Precalculus}
\newcommand{\MeetingLength}{60 minutes}
\newcommand{\UnitNumberName}{Unit 1: Functions and Change \quad}
\newcommand{\LessonNumberName}{Lesson 1.2: Function Operations}

\begin{document}
\begin{center}
    {\Large\bfseries \CourseName \\
    {\small\bfseries \UnitNumberName \LessonNumberName}}
\end{center}

\begin{tcolorbox}[colback=lilac, colframe=plum, boxrule=0.9pt, arc=2mm,
    left=2mm, right=2mm, top=1.5mm, bottom=1.5mm]
    \textbf{Primary Objective:} Students will add, subtract, multiply, and divide
    two functions --- evaluating a combination at a point from a table and
    building it as a formula --- and will state the \textit{domain} of the result
    from the original functions: the overlap of both parents' domains, with every
    input that makes the bottom function zero thrown out of a quotient.
    \hfill\textit{\small (Foundations ramp --- no CED topic)}
\end{tcolorbox}

\boxguard[22]
\begin{skillbox}[Lesson Flow --- Gradual Release (60 minutes)]{lilac}
\small
\renewcommand{\arraystretch}{1.25}
\begin{tabularx}{\linewidth}{@{} l c X @{}}
\textbf{Phase} & \textbf{Min} & \textbf{What students are doing, and where it lives} \\
\midrule
Warm-Up                       & 5  & Silent, individual spiral review --- warm-up sheet. \\
Hook                          & 3  & Whole-class discussion --- hook box atop the guided notes. \\
\textbf{I Do} --- model       & 10 & Watch and annotate while the teacher thinks aloud --- the COMBINING and AT A POINT rows: definitions, the machine diagram, problems 1 and 3. \\
\textbf{We Do} --- guided     & 21 & Fill the notes together, every problem cold-called --- problems 2, 4--6, then the THE FORMULA row (7--10) and the DOMAIN row (11--14; the teacher works 11, the class works the crux, 14). \\
\textbf{You Do} --- independent & 13 & Work alone while the teacher circulates --- the You Do row, problems 15--18. \\
Debrief                       & 5  & Whole-class share-out and the headline. \\
Homework Launch               & 3  & Start homework problems 1 and 5 alone; the teacher reads 5(b) over shoulders. \\
\end{tabularx}

\smallskip
\textbf{Protect the DOMAIN row.} Three of the four instruction rows are
procedural and will go faster than you expect; the fourth is the entire lesson.
If the clock slips, cut \textbf{problem 10} (the product as a formula) and then
\textbf{problem 9}, both of which the homework re-asks in Part B --- and protect
problems \textbf{13 and 14}, the cancel case. A class that never sees the two
answers disagree has not met the misconception at all. Never cut the Homework
Launch to zero: with no exit ticket it is the last formative read of the period.
\end{skillbox}

\renewcommand{\arraystretch}{1.6}
\begin{skillbox}[Priority Ideas \& Skills]{goldbox}
\begin{tabularx}{\linewidth}{@{}X X@{}}
\textbf{Priority Skills} &
\textbf{Key Understandings} \\
\midrule
\textbf{AP Skill 1.A} --- Identify and describe a function built from other functions. &
\begin{itemize}[leftmargin=1.2em,itemsep=2pt,topsep=0pt]
  \item $(f+g)(x) = f(x) + g(x)$, and likewise for $-$, $\times$, $\div$: every
        operation is arithmetic done to the two \textit{outputs}.
  \item Both functions must accept the same input before anything can be
        combined.
\end{itemize} \\
\textbf{AP Skill 2.B} --- Construct equivalent representations of a combination. &
\begin{itemize}[leftmargin=1.2em,itemsep=2pt,topsep=0pt]
  \item At a point, two table lookups are enough --- no formula is needed.
  \item As a formula, each function's whole rule goes in its own parentheses
        before anything is distributed.
  \item \textbf{Secondary trap:} $(f-g)$ and $(g-f)$ are different functions,
        exact opposites of each other. The same is true of $f/g$ and $g/f$.
\end{itemize} \\
\textbf{AP Skill 3.B} --- Justify a conclusion about a domain with appropriate reasoning. &
\begin{itemize}[leftmargin=1.2em,itemsep=2pt,topsep=0pt]
  \item The domain of a combination is the \textbf{overlap} of the two parents'
        domains --- both must say yes.
  \item For a quotient, additionally solve $g(x) = 0$ and throw those inputs out.
  \item \textbf{Target misconception:} students read the domain of a combination
        off the \textit{simplified} formula rather than off the \textit{original}
        functions. With $f(x) = x^2 - 9$ and $g(x) = x - 3$, the quotient
        simplifies to $x + 3$, which looks defined everywhere --- but $g(3) = 0$,
        so $x = 3$ is still excluded. The two readings \emph{disagree}, and that
        disagreement is the lesson.
\end{itemize} \\
\end{tabularx}
\end{skillbox}

\begin{skillbox}[{Vocabulary, Concepts, \& Theorems}]{greenbox}
\renewcommand{\arraystretch}{1.4}
\begin{tabularx}{\linewidth}{@{} >{\leavevmode\bfseries\color{plum}}p{0.24\linewidth} X @{}}
Sum of two functions & \TallMath{(f+g)(x) = f(x) + g(x)} --- add the two outputs at the same input. \\
Difference of two functions & \TallMath{(f-g)(x) = f(x) - g(x)}; order matters, and $g-f$ is its opposite. \\
Product of two functions & \TallMath{(fg)(x) = f(x) \cdot g(x)} --- multiply the two outputs. \\
Quotient of two functions & \TallMath{\left(\frac{f}{g}\right)(x) = \frac{f(x)}{g(x)}}, defined only where $g(x) \ne 0$. \\
Domain of a combination & The overlap of both parents' domains, with every input making the bottom function zero removed as well. \\
Profit function & A difference built in context: $P = R - C$, revenue minus cost. \\
\end{tabularx}
\end{skillbox}

\begin{skillbox}[Activate Prior Knowledge \& Spiral Review (5 min)]{lilac}
    Please see attached warm-up (spiral review).
    Skills reviewed: evaluating a function at a negative input,
    subtracting a whole expression by distributing the minus sign,
    multiplying two binomials,
    and solving an equation of the form $x - 3 = 0$ ---
    the exact move that finds the inputs a quotient must throw out.
\end{skillbox}

\begin{skillbox}[Hook (3 min)]{lilac}
    \textbf{The owner wants one number.} A bakery bakes $x$ loaves in a day and
    tracks two costs separately: flour, $F(x) = 0.80x$ dollars, and labor,
    $L(x) = 25 + 1.20x$ dollars.
    \begin{enumerate}[itemsep=2pt]
        \item Put both rules on the board and say the complaint out loud: ``I do
              not want two numbers every morning --- I want \textit{one}.'' Ask
              what the owner should do.
        \item Take ``add them'' and then press: add \textit{what}, exactly? The
              rules, or the outputs? Land on
              $T(x) = F(x) + L(x) = 25 + 2x$.
        \item Ask the question that licenses the whole lesson: why is it legal to
              add these two at all? Because both take the \emph{same} input,
              $x$ loaves.
    \end{enumerate}
    \textit{Discussion prompt:} we have spent a lesson running one machine at a
    time. Today two machines run on the same input and their outputs are
    combined into a third function.
\end{skillbox}

\begin{skillbox}[Lesson --- I Do / We Do / You Do]{lilac}
\begin{multicols}{2}

\textbf{I DO --- Part 1: Two Machines, One Answer} (10 min)

\textit{Notes rows COMBINING and AT A POINT: read each definition, mark up the
display, work problems 1 and 3. Teacher works; students watch and annotate.
Nothing is cold-called in this phase.}
\begin{itemize}
  \item Read the four definitions aloud, slowly, and say the same sentence after
        each one: \textit{the operation happens to the outputs}.
  \item Mark up the machine diagram yourself. Fill the two label boxes --- both
        machines are fed the same \textbf{input $x$}; what meets at the combiner
        is the two \textbf{outputs}. Trace the input arrow splitting in two with
        your finger before you say a word about arithmetic.
  \item Work \textbf{problem 1} aloud: writing $(f+g)(x) = f(x) + g(x)$ is the
        answer, not a restatement of the question.
  \item Read the AT A POINT sentences and work \textbf{problem 3} from the table:
        $(f+g)(4) = -3 + 16 = 13$. Narrate the two lookups \textit{before} the
        arithmetic, so the class sees there is no formula anywhere.
\end{itemize}

\textbf{WE DO --- Part 2: At a Point, and the Order Trap} (7 min)

\textit{Problems 2 and 4--6. Class fills the blanks together; cold-call every
one.}
\begin{itemize}
  \item \textbf{Problem 2} first: the bakery total, now written as $(F+L)(x)$.
        The hook becomes notation.
  \item Fill the $(f+g)$ row of the table together --- four values, cold-called,
        no formula.
  \item \textbf{Problems 5 and 6} are the order trap as a pair: $(f-g)(2) = -3$
        and $(g-f)(2) = 3$. Put both on the board at once and make the class say
        what changed. Do not let ``they are the same but backwards'' stand ---
        push to ``opposites.''
\end{itemize}

\columnbreak

\textbf{WE DO --- Part 3: Building the Rule} (6 min)

\textit{Notes row THE FORMULA. Class fills the blanks together; cold-call every
one. Teacher works problem 7; the class works 8--10.}
\begin{itemize}
  \item Read the three steps, then insist on the parentheses in step 2 before
        anybody simplifies. The display's gold note is the whole risk:
        $-(x-3)$ is $-x + 3$.
  \item \textbf{Problem 8} is where the minus sign gets distributed; make one
        student read the second line aloud, sign by sign.
  \item \textbf{Problem 9} is the order trap again, now at the formula level:
        $-x - 4$ is the exact opposite of problem 8's $x + 4$. Say it back.
  \item \textbf{Problem 10} is the product. Cut this one first if the clock is
        tight --- homework problem 4 asks for it again.
\end{itemize}

\textbf{WE DO --- Part 4: Who Is Allowed In} (8 min)

\textit{Notes row DOMAIN --- the spine of the lesson. Class fills the blanks
together; cold-call every one. Teacher works problem 11; the class works
12--14.}
\begin{itemize}
  \item Do the number lines before any algebra: shade the two parent domains,
        then the overlap, and fill the label box --- the overlap is the domain of
        \textbf{$f+g$}. Both parents must say yes.
  \item Then the quotient's extra rule: solve $g(x) = 0$ and throw those inputs
        out. This is warm-up item 4, arriving with a job.
  \item Mark up the cancel case on the right of the display. $g(3) = 0$, so the
        quotient never had a value at 3 --- and the open circle is still there
        after the cancelling. Fill the second label box: \textbf{undefined}. Give
        one minute for the ``in your own words'' line.
  \item Work \textbf{problem 11}. \textbf{Problem 13} does the cancelling;
        \textbf{problem 14 is the crux} --- ask for the \textit{reason} 3 is
        excluded, never just the answer. Say the headline here for the first
        time.
\end{itemize}

\textbf{YOU DO --- Part 5: On Your Own} (13 min)

\textit{Notes row ON YOUR OWN, problems 15--18 --- the whole independent phase.
Students work alone; circulate and say nothing a neighbor could say instead.}
\begin{itemize}
  \item \textbf{Problem 15} is the formula skill unassisted; \textbf{16} is the
        order trap met alone.
  \item \textbf{Problem 17 is the crux transferred} --- $x^2 - 25$ over $x - 5$,
        with no display to lean on and nobody saying ``check the bottom.''
  \item \textbf{Problem 18} is the finisher's question: build $P = R - C$ for the
        food truck, evaluate it, and give a domain that a truck could actually
        honor. It returns as homework problem 7.
\end{itemize}
\end{multicols}
\end{skillbox}

\begin{skillbox}[Active Monitoring]{redbox}
Circulate during the \textbf{You Do} phase --- the thirteen minutes where students
are working problems 15--18 without you --- and check. The crux is
\textbf{problem 14} (the domain of $f/g$ is all reals except 3, even though the
quotient simplifies to $x+3$); its You Do transfer is \textbf{problem 17}
($p/q$ with $q(5) = 0$, no display on the page).
\begin{itemize}
  \item \textbf{Common error (the crux):} at problem 17, simplifying to $x + 5$
        and then writing ``all real numbers.'' Push: ``What is $q(5)$? Was 5 ever
        an input you were allowed to use?'' Point back at the open circle they
        annotated in Part 4.
  \item \textbf{Common error:} at problem 16, giving the same answer twice, or
        computing $(g-f)$ when $(f-g)$ was asked. Ask which function is written
        \emph{first}, then which one is being subtracted.
  \item \textbf{Common error:} dropping the parentheses in problem 18, so
        $P(x) = 9x - 4x + 120$. Make them write $(4x + 120)$ inside parentheses
        before subtracting anything.
  \item \textbf{Common error:} an unrestricted domain on problem 18 --- a truck
        cannot sell $-4$ or $7.5$ plates, and 60 is the cap. Ask: ``Would the
        truck accept that number of plates?''
  \item \textbf{Watch for:} a student who computes $(fg)(x)$ in problem 15 by
        multiplying only the first terms and only the last. Have them say
        ``every term times every term'' as they write.
\end{itemize}
\end{skillbox}

\begin{skillbox}[Differentiation --- During You Do (13 min)]{redbox}
There is no separate activity sheet: students work the notes' You Do row
(problems 15--18) alone while you circulate. Differentiate by \textit{where you
stand}, not by handing out different paper.

\begin{itemize}
  \item \textbf{Support.} Sit with a struggling student on \textbf{problem 15}
        and make them write $(3x - 2) + (x + 7)$ with both sets of parentheses
        before a single like term is combined. A student who cannot get through
        15 will not get through 17, and the parentheses are the whole fix.
  \item \textbf{On level.} \textbf{Problems 16 and 17} are what the class should
        clear unaided. Do not answer 17. Ask one question --- ``what is $q(5)$?''
        --- and walk away before they answer it.
  \item \textbf{Extend.} \textbf{Problem 18} is the finisher's question: it asks
        for a context domain, which no earlier problem did. For a student who
        finishes it, hand over this one: \textit{if the truck's cost were
        $C(x) = 4x + 120$ and its revenue were $R(x) = 9x$, what does
        $\left(\frac{C}{R}\right)(x)$ mean, and why must $x = 0$ be thrown out of
        its domain?} That answer presents in the Debrief.
\end{itemize}
\end{skillbox}

\boxguard[20]
\begin{skillbox}[Debrief (5 min)]{redbox}
Whole class, teacher-led, packets still open. Three moves, in order:
\begin{enumerate}[itemsep=3pt]
  \item \textbf{Share out the You Do} (3 min) --- one answer per problem, not a
        full review. Problem 15: the sum and the product. Problem 16: both
        differences, said aloud as opposites. \textbf{Problem 17: the simplified
        form and the domain, in that order and as two separate facts} --- this is
        the one every student must hear even if they never reached it. Problem
        18: $P(x) = 5x - 120$, $P(40) = 80$, and a domain a food truck could
        honor. If a student answered the extension question, they present it.
  \item \textbf{Name the headline} (1 min) --- say it, then make the class say it
        back: \textit{the domain of a combination comes from the original
        functions, never from the simplified formula. Simplifying can hide a
        forbidden input; it can never make it legal.} Point at the open circle on
        the notes display while you say it.
  \item \textbf{Point forward} (1 min) --- 1.3 stops feeding both machines the
        same input. The output of one becomes the \textit{input} of the next.
\end{enumerate}
\textbf{Do not} re-teach the formula work here, take new questions, or let the
homework start early. If the launch shows the domain did not land, it is the 1.3
warm-up's job, not this minute's.
\end{skillbox}

\begin{skillbox}[Homework Launch (3 min)]{redbox}
\textbf{There is no exit ticket.} The formative read comes twice: circulating the
You Do, and again here. Write the due date on the board, have students copy it
into the cover's Homework row, then start \textbf{homework problems 1 and 5}
alone and in silence (3 minutes). Circulate straight to the students you flagged
during You Do.
\begin{itemize}[itemsep=2pt]
  \item \textbf{Problem 1} is the quick read: two table lookups per part. A
        student still hunting for a formula has not heard Part 2.
  \item \textbf{Problem 5(b) is the diagnostic}: the domain of $f/g$ for
        $f(x) = x^2 - 16$ and $g(x) = x - 4$. The answer is all real numbers
        except 4. A student who writes ``all real numbers'' read it off the
        simplified $x + 4$ --- the exact misconception --- and a student who
        writes the right domain but cannot name $g(4) = 0$ as the reason has
        copied a pattern rather than understood it.
\end{itemize}
\textbf{Sort what you see into three piles} as you walk: \textit{domain from the
originals} (nothing to do), \textit{simplifies first, then reads} (a 1.3 warm-up
item on excluded inputs), and \textit{still losing the minus sign} (a one-minute
check-in at the door of 1.3).
\end{skillbox}

\begin{skillbox}[Reinforcement \& Extension]{goldbox}
\textbf{Homework --- printed, graded (2 pages).} Last in the packet; due on the
date the teacher sets. Every context is new. \textit{Part A} (1--3) combines at a
point from a table, repeats the order trap at $x = 3$, and builds a print shop's
total cost as a sum. \textit{Part B} (4--6) builds formulas and their domains:
all four operations on one pair, then \textbf{problem 5, the diagnostic} --- the
quotient that cancels --- then a pair whose sum and quotient have different
domains. \textit{Part C} (7--8) is context: a lawn-care profit $P = R - C$ with a
context domain, and a car wash's average cost per car, where $x = 0$ must go.
\textit{Part D} (9--10) puts a classmate's wrong domain on the page to be
corrected, and problem 10 asks for the headline in writing.

\textbf{AP Practice --- extra credit (2 pages).} Optional, before the homework in
the packet. Section I is five AP-style multiple-choice items (four options) whose
distractors are the lesson's errors: adding or reversing when a difference was
asked (1), answering $(g-f)$ when $(f-g)$ was asked (2), reading the domain off
the \emph{simplified} quotient (3 --- the crux), adding when a product was asked
or dropping the middle term (4), and forgetting that a parent's own excluded
input is excluded from the sum as well (5). Section II is one free-response
question on a bike shop's revenue and cost table: interpret a difference in
context (a), compare $R-C$ with $C-R$ (b), locate where profit turns positive
(c), explain why a quotient dies at $x = 0$ (d), and build a profit formula with
its context domain (e) --- the finisher's question.

\textbf{Preview:} Next lesson (1.3, Function Composition and Decomposition)
chains the two machines instead of combining them: the output of one becomes the
\textit{input} of the next, $f(g(x))$, and a complicated rule is decomposed back
into the two simpler machines it was built from. The homework's closing box
points back at the print shop in problem 3 to set it up.
\end{skillbox}

\begin{teachernote}[Warm-Up]
Five minutes, silent start, no calculators. Item 2 is the one to watch: a student
who writes $4x + 9 - x - 6$ has lost the minus sign, and will lose it again
inside $(f-g)(x)$ within twenty minutes --- note names now. Item 4 looks like the
easiest thing on the page and is the most important: solving $x - 3 = 0$ is
\emph{exactly} what finding the excluded input of a quotient asks for, and part
(b) makes the point that the answer is not always the number you see written in
the expression. If item 4 is slow, say out loud that you will need it again in
twenty minutes, and mean it.
\end{teachernote}

\begin{teachernote}[Guided Notes]
The notes are one Main Ideas / Notes table, four instruction rows and a You Do
row; in each instruction row you read the printed definition, mark up the
display, and work the first problem of the grid (1, 3, 7, 11), and the class
works the rest with the pen in their hand, every problem cold-called. The first
three rows are procedural and will move quickly --- resist the urge to let them
expand, because the fourth row is the lesson.

COMBINING and AT A POINT are one block: the point of the table in row 2 is that
no formula exists anywhere on that page and the combination still has a value.
Problems 5 and 6 are a deliberate pair; put both answers on the board at once,
because ``opposites'' is not obvious to a student who has just been told the two
expressions ``mean the same thing.'' In THE FORMULA row, the only real risk is
the distributed minus sign, and the gold box in the display is there to be
pointed at.

DOMAIN is the spine, and it must not be rushed. Do the two number lines first ---
overlap is a visual idea and the algebra comes easier after it. Then the cancel
case. Do not simplify and then ask about the domain; ask about the domain
\emph{first}, get ``all reals except 3'' from the original $g$, and only then
cancel, so the class watches the evidence disappear from the page while the
exclusion stays. \textbf{Problem 14 is the crux} and it asks for a reason, not a
number: do not accept ``because there's a hole.'' Make them say $g(3) = 0$.

The You Do row (15--18) is the entire independent phase and carries thirteen
minutes. Problem 17 is the crux with the scaffolding removed; problem 18 is the
first time anyone has been asked for a \emph{context} domain, which is genuinely
harder than it looks. If the clock slips, cut problem 10 from row 3 rather than
anything in row 4 or 5.
\end{teachernote}

\begin{teachernote}[Debrief]
Five minutes, and they are at risk because the DOMAIN row will want them. Borrow
from problem 10 instead --- the product formula is the most replaceable thing in
the period, and homework problem 4 asks for it again. Do not borrow from the
three-minute Homework Launch either: with no exit ticket, it is the last look you
get before the work leaves the room. The share-out is one answer per You Do
problem specifically so that a student who stalled on 15 still hears problem 17's
two-part answer --- the simplified form \emph{and} the domain, stated separately.
That pairing is the shape of homework problem 5, which they will otherwise meet
alone at home.
\end{teachernote}

\begin{teachernote}[AP Practice]
Extra credit, optional, and not a substitute for the homework --- say so when you
hand out the packet. Score it however extra credit works in your gradebook; the
cover gives it a $+$ cell outside the total. Multiple-choice answers are
\textbf{B, D, B, C, A}. Item 3 is the one to glance at: a student who circles (A)
``all real numbers'' cancelled to $x + 7$ and read the domain off it --- that is
the lesson's misconception in its purest form, and it means problem 14 of the
notes did not land. A student who circles (D) has excluded $-7$ as well, which is
the \emph{numerator's} zero and harmless. On item 2, (A) is $(g-f)(2)$ and is
worth counting: if more than a few take it, spend a warm-up item on order in 1.3.
In the free response, part (c) needs values, not a sentence --- require both
$P(5)$ and $P(10)$ --- and part (e) is the deepest: accept any domain that
restricts to whole bikes with $0 \le x \le 20$.
\end{teachernote}

\begin{teachernote}[Homework]
Graded, two pages, due on the date you set; students start problems 1 and 5 in
the Homework Launch. Grade \textbf{problems 5 and 9} for accuracy and the rest
for completion and shown work if time is short: 5(b) and 9 are the same question
asked from two sides --- state the domain yourself, and catch somebody else
stating it wrongly --- and together they predict 1.3 better than any of the
formula items do. Problem 10 is the headline in writing; read it for the words
\emph{overlap} and \emph{zero}, not for polish. On problem 6, a student who gives
the sum and the quotient the same domain has forgotten the quotient's extra rule;
if that is common, open 1.3 with it. Problem 7(c) is the context domain and the
most commonly skipped item on the page --- a bare ``all real numbers'' there is
worth one sentence of feedback, because 1.3's compositions will ask for context
domains repeatedly.
\end{teachernote}
\end{document}
